\chapter*{Liste des acronymes}
\addcontentsline{toc}{section}{\textbf{Liste des acronymes} }
  
\begin{spacing}{1.05}

\begin{description}
  \item[MP] \hfill \\
  \textbf{Matière Première} : établissement public de santé regroupant les hôpitaux de Castres et de Mazamet, porteur du programme RI2S.

  \item[MPU] \hfill \\
  \textbf{Matière Première Usine} : programme d’expérimentation visant à tester des solutions numériques pour la santé des seniors en zone rurale.

  \item[MPF] \hfill \\
  \textbf{Matière Première Fournisseur} : interface logicielle permettant à différentes applications ou composants de communiquer entre eux via des appels standards.

  \item[PLM] \hfill \\
  \textbf{Project Life Management} : dispositifs de coordination entre professionnels de santé à l’échelle d’un territoire, pour améliorer la prise en charge des patients.

  \item[ETL] \hfill \\
  \textbf{Extract Transform Load} : langage de feuilles de style utilisé pour décrire la présentation et la mise en forme des pages HTML.

  \item[BI] \hfill \\
  \textbf{Business Intelligence} : langage de balisage standard utilisé pour structurer le contenu d’une page web.

  \item[BU] \hfill \\
  \textbf{Business Unit} : réglementation européenne encadrant le traitement et la circulation des données à caractère personnel.

  \item[DCPC] \hfill \\
  \textbf{DermoCosmétique Personnal Care} : langage de modélisation graphique utilisé pour représenter visuellement la structure, les interactions et le comportement d’un système logiciel.

  \item[UX/UI] \hfill \\
  \textbf{User Experience / User Interface} : ensemble des méthodes visant à optimiser l’ergonomie, l’utilisabilité et l’interface visuelle d’un produit numérique.
\end{description}
\end{spacing}